\subsection{Fixed operators and extension-dependent quantum grading}
Clearing the normalization denominators in \fid{F26} removes
the obstruction for an individual word. It does not produce
one finite Frobenius space supporting all words under the same
extension-dependent substitution.

\begin{finding}[\fid{F42}: a common finite spectrum cannot encode all words]\label{f:commonwords}
On the actual rank-three tensor module of \fid{F21}, put
$D=v^2+2+v^{-2}$. For every $k\ge1$,
\[
 \Tr(P_1^k)=40D^k.
\]
Fix $q>1$ and set $D_n=q^n+2+q^{-n}$.
There is no finite-dimensional complex vector space $V$,
fixed invertible $F\in\End(V)$, and endomorphisms $T_k$ for
an unbounded set of positive integers $k$, such that
\[
 \Tr(T_kF^n)=40D_n^k
\]
for all sufficiently large $n$. The threshold may depend on $k$.
The assertion permits noncommuting $T_k,F$ and Jordan blocks,
and remains valid for finite Euler supertraces. It also remains
valid after multiplying each right-hand side by a nonzero scalar
$a_k$ and a single nonzero exponential factor $\beta_k^n$.

When $q$ is a prime power, the entire actual tensor family at the single
parameter $v_0=q^{1/2}$ has an exact action-compatible realization
by finite Weil objects with identity Frobenius. Its normalized
generators are genuine idempotent endomorphisms. Moreover each
individual sequence $40D_n^k$ has a positive finite Tate
coefficient object, with that object allowed to depend on $k$.
These constructions do not supply the opposite-sign path
comparison required in \fid{F24}.
\end{finding}

\begin{proof}
Use the consistent normalization of \fid{F19}: $c_i^2=\delta c_i$,
$\delta=v+v^{-1}$, and
\[
 P_i=c_i\otimes c_i+(c_i-\delta)\otimes(c_i-\delta).
\]
The two terms are orthogonal and their squares are $D$ times
themselves, so $P_i^2=DP_i$. This uses the symbolic definition
in GCT IV, equation~(93), with the section-3.1 convention
\cite[\S3.1 and Definition~11.1, equation~(93)]{IV},
rather than silently using the inconsistent section-11
assignments already recorded in the source audit.
In the decomposition of \fid{F21}, $P_1$ acts by $D$ on
the scalar positive block of multiplicity 20 and by zero on
the scalar negative block of multiplicity 4. Each of the
two-dimensional blocks has rank one for $P_1$ and trace $D$,
with multiplicities 16 and 4. Hence
\[
 \Tr(P_1^k)=40D^k\qquad(k\ge1).
\]
Both substitutions $v=q^{n/2}$ and $v=q^{-n/2}$ give
\begin{equation}\label{eq:wordpowers}
 40D_n^k
   =40\sum_{j=0}^{2k}\binom{2k}{j}q^{(k-j)n}.
\end{equation}
The $2k+1$ roots in this sum are distinct, and every coefficient
is nonzero.

For fixed $F,T_k$, the generating function beginning at $n=0$ is
\[
 \sum_{n\ge0}\Tr(T_kF^n)T^n
   =\Tr\bigl(T_k(1-TF)^{-1}\bigr).
\]
All its poles are among the inverses of the finitely many
eigenvalues of $F$. A Jordan block can increase a pole order
but cannot add a pole location. Removing finitely many initial
terms changes only a polynomial. On the other hand the
generating function of \eqref{eq:wordpowers} is
\[
 40\sum_{j=0}^{2k}\binom{2k}{j}
      \frac{1}{1-q^{k-j}T}.
\]
Every one of its $2k+1$ poles is genuine: only one summand has
that pole, and its coefficient is nonzero. Thus a common $F$
would require at least $2k+1$ distinct eigenvalues for unbounded
$k$, which is impossible.
For a finite graded object use the union of the even and odd
spectra; for a bounded complex use its finite-dimensional
cohomology. Multiplication by $a_k\beta_k^n$ preserves the
number of distinct nonzero roots, now
$\beta_kq^{k-j}$, proving the normalization variant.
This proof even permits unrelated $T_k$; requiring
$T_k=P^k$ or additional tensor relations cannot evade it.

We give the positive fixed-parameter construction with its
action comparison stated explicitly. Let $q$ be a prime power,
$\ell\nmid q$, and let a finite extension $E/\Q_\ell$
contain $v_0=q^{1/2}$.
Write
\[
 M_{r,q}=(V_q\otimes W_q)^{\otimes r}
\]
for the actual tensor module with the two ordinary Hecke
actions at $v_0$, and let $A_{r,q}$ be the image algebra of its
specified nonstandard generators. An action comparison consists
of finite Weil objects $H_r$, fiber identifications
$\omega(H_r)\simeq M_{r,q}$, and right $A_{r,q}$-actions
by Weil endomorphisms such that:
\begin{enumerate}
\item every named $P_i$ is intertwined by the fiber
identification;
\item the tensor identifications $H_r\otimes H_s\simeq H_{r+s}$
intertwine the local generators on disjoint consecutive blocks,
with the usual tensor associativity coherence.
\end{enumerate}
Extension of the finite field raises Frobenius to its $n$th
power while keeping $v_0$ and these operators fixed.

Give $V_q,W_q$ identity Frobenius and use the induced identity
Frobenius on their tensor powers. Every actual matrix is then
a Weil endomorphism. The local Hecke formula in
GCT IV, \S6.1 \cite[\S6.1, p.~33]{IV}, gives the actual operators, and locality proves
the disjoint-block compatibility. All their algebra relations
are unchanged. This proves the stated action comparison for
the whole actual image-algebra family.
For fixed dimensions two these are the two-row quotient actions;
no faithfulness assertion for the full abstract algebra at
every rank is required. The quantum Hecke operators are extra
morphisms, not an identification with the symmetry maps of the
ordinary Weil category.

Since $D_0=(v_0+v_0^{-1})^2\ne0$, scalar division defines
$e_i=P_i/D_0$ and gives $e_i^2=e_i$. Their images are genuine
Weil subobjects. This divides by a coefficient-field scalar;
it does not tensor-invert the object of \fid{F38}.
The resulting extension trace is
\[
 \Tr\bigl(w(P_i(v_0))F_r^n\bigr),
\]
with fixed operators. For identity Frobenius it is constant
in $n$, not the formal word trace retuned at $q^{n/2}$.
In particular the normalized mixed-word trace of \fid{F26}
is the constant $24+16/D_0^2$. This is an endomorphism trace,
not automatically the Euler multiplicity of an object;
$e_1e_2$ need not be idempotent.

There is also a valid coefficient-object construction for
\eqref{eq:wordpowers}. On a point let
\[
 \mathcal D=E(-1)\oplus E^{\oplus2}\oplus E(1),\qquad
 \mathcal C_k=E^{\oplus40}\otimes\mathcal D^{\otimes k}.
\]
Geometric Frobenius has trace $D_n$ on $\mathcal D$ and
$40D_n^k$ on $\mathcal C_k$ at every extension.
The canonical tensor identification
$\mathcal C_{k+1}\simeq\mathcal D\otimes\mathcal C_k$
realizes multiplication by this coefficient.
Its dimension is $40\cdot4^k$ and its weights range from
$-2k$ to $2k$, so positivity here does not assert purity.

The one-generator quadratic relation itself has a coherent
functor model. On the product of 64 finite Weil-representation
categories, let $\mathcal E_{40}$ project onto 40 chosen
components and set
$\mathcal P=\mathcal D\otimes\mathcal E_{40}$. Then
\[
 \mathcal P\circ\mathcal P
    \simeq\mathcal D\otimes\mathcal P
\]
by tensor associativity and
$\mathcal E_{40}^2=\mathcal E_{40}$, compatibly with further
powers. This models the split action of one generator and its
quadratic relation. It does not identify a simultaneous model
for the second generator, the original integral tensor lattice,
or the designated nonstandard bases.

Finally the action comparison alone cannot imply the main
sign hypothesis. Take a constant rank-$N$ sheaf with identity
Frobenius on affine $d$-space, with any of the preceding
fixed-parameter actions on its fibers. Its tensor and named
operator relations remain exact. Its pointwise trace is $N$,
so the pointwise hypothesis of \fid{F24} holds with $a=0$
and $B=N$, while its total compact trace is
\[
 M_n=Nq^{dn}.
\]
For $d>0$ the visible root $q^d$ exceeds $q^{d/2}$.
The proved \fid{F24} therefore rules out the required
opposite-sign path comparison with sublinear boundary budget
on an unbounded even set for this model. This proves the
claimed insufficiency with an explicit action-compatible
example.
\end{proof}

Thus a common finite operator realization, a fixed-parameter
action, and word-dependent coefficient objects impose different
comparison requirements. The latter objects may grow with the
word; the common-space obstruction does not exclude them.
For the main route one must additionally identify the actual
path complex and its tensor or idempotent construction by a
Weil-compatible comparison, with the required odd shift and a
controlled cone. Equality of fixed-operator traces alone does
not provide that comparison. Nor does fixed-$q$ numerical
positivity imply coefficientwise positivity in a formal
variable: $x^2-x+1$ is positive for every real $x$ but has a
negative coefficient.

% Separate paragraph for Part III.
\paragraph{\fid{F42} in Grothendieck's language.}
The finite-resolvent obstruction concerns one finite realization
and its fixed Frobenius, whereas the Tate coefficient objects
use tensor powers with growing spectra. The fixed-parameter
construction is an exact action in the category of point Weil
objects, with named generators and tensor compatibility; it is
not an extension-wise quantum grading comparison. Scalar
division, an idempotent image, and inversion of a tensor object
are distinct operations here. The remaining path comparison
must preserve the actual operations and control its boundary
complex; no such comparison is attributed to an archive leaf.
